\usepackage{graphicx} % Required for inserting images
\usepackage{float}
\usepackage{amsmath}

\usepackage{tocloft}
\usepackage{titlesec}
\usepackage{lmodern}

\usepackage{multirow}

\usepackage{array}
\newcolumntype{C}[1]{>{\centering\arraybackslash}p{#1}}
\newcolumntype{R}[1]{>{\raggedleft\arraybackslash}p{#1}}

\usepackage[hidelinks]{hyperref}

\usepackage{enumitem}
\usepackage{tikz}
\usetikzlibrary{shapes, arrows, positioning}

% ---- Paramètres de marges ----
\usepackage{geometry}
\geometry{
    a4paper,
    left=2.5cm,
    right=2.5cm,
    top=2.5cm,
    bottom=2.5cm
}

% ---- Paramètres des tableaux ----

\usepackage{longtable}
\usepackage[table]{xcolor}
\setlength{\tabcolsep}{10pt}
\renewcommand{\arraystretch}{1.5} 

% ---- Packages de mise en page ----
\usepackage[utf8]{inputenc}
\usepackage[T1]{fontenc}
\usepackage[french]{babel}
\addto\captionsfrench{
  \renewcommand{\tablename}{Tableau}
}

% ---- Numerotation des chapitres ----
\renewcommand{\thechapter}{\arabic{chapter}}
\renewcommand{\thesection}{\thechapter.\arabic{section}}

% ---- En-têtes et pieds de pages ----
\usepackage{fancyhdr}
\setlength{\headheight}{16pt}

\pagestyle{fancy}
\fancyhf{}
\fancyhead[C]{\leftmark}
\fancyfoot[C]{\thepage}
\renewcommand{\headrulewidth}{0.4pt}

\fancypagestyle{plain}{
  \fancyhf{}
  \fancyhead[C]{\leftmark}
  \fancyfoot[C]{\thepage}
  \renewcommand{\headrulewidth}{0.4pt}
}

% ---- Paragraphe et indentation ----
\usepackage{setspace}
\setstretch{1.5}            % Interligne 1.5
\setlength{\parindent}{0pt} % Sans retrait
\setlength{\parskip}{1em}   % Espacement entre les paragraphes
\hyphenpenalty=10000
\exhyphenpenalty=10000

% ---- Table des matieres ----

% Espace avant chaque chapitre
\setlength{\cftbeforechapskip}{8pt}
% Espace avant sections
\setlength{\cftbeforesecskip}{4pt}
% Espace avant sous-sections
\setlength{\cftbeforesubsecskip}{2pt}
% Espace apres nemuro
\setlength{\cftfignumwidth}{3em}
% 1. Augmenter la largeur pour "Chapitre VII"
\setlength{\cftchapnumwidth}{2em}
% 1. Augmenter la largeur pour "Chapitre VII"
\setlength{\cftsubsecnumwidth}{3.5em}

% ---- Centrage des titres auto-générés (tocloft) ----

% Pour la Table des Matières
\renewcommand{\cfttoctitlefont}{\hfill\Huge\bfseries}
\renewcommand{\cftaftertoctitle}{\hfill\markboth{TABLE DES MATIÈRES}{TABLE DES MATIÈRES}}

% Pour la Liste des Figures
\renewcommand{\cftloftitlefont}{\hfill\Huge\bfseries}
\renewcommand{\cftafterloftitle}{\hfill\markboth{LISTE DES FIGURES}{LISTE DES FIGURES}}

% Pour la Liste des Tableaux
\renewcommand{\cftlottitlefont}{\hfill\Huge\bfseries}
\renewcommand{\cftafterlottitle}{\hfill\markboth{LISTE DES TABLEAUX}{LISTE DES TABLEAUX}}

% ---- Style chapitre ----
\titleformat{\chapter}[hang]
{\normalfont\large\bfseries\centering} % \large correspond à la taille 14pt
{}{0em}{}

\makeatletter
\def\@makechapterhead#1{%
  \thispagestyle{empty}
  \vspace*{\fill}
  \vspace*{-50pt}
  {\centering
    \normalfont\large\bfseries
    Chapitre \thechapter\par
    \vspace{10pt}
    \large #1\par
  }
  \vspace*{\fill}
  \newpage
}
\makeatother

\makeatletter
\def\@makeschapterhead#1{%
  \thispagestyle{fancy}
  \markboth{\MakeUppercase{#1}}{\MakeUppercase{#1}}
  
  \vspace*{-15pt} 
  
  {\centering
    \normalfont\large\bfseries
    #1\par
  }
  
  \vspace{25pt} 
}
\makeatother

% ---- Style itemize ----

\usepackage{enumitem}

\setlist[itemize]{
    label=$\cdot$,
    itemsep=5pt,    % Espace entre les items
    parsep=0pt,     % Espace entre les paragraphes à l'intérieur d'un item
    leftmargin=1cm
}

% ---- Style Tableaux ----

\usepackage[table]{xcolor}

% \definecolor{bleuEntete}{RGB}{31, 73, 125} % Un bleu marine profond et pro
\definecolor{bleuEntete}{RGB}{79, 129, 189} % Un bleu équilibré et moderne

% Définition d'un bleu personnalisé (optionnel, plus pro que le "blue" par défaut)
\definecolor{bleuTresClair}{RGB}{240, 245, 255} 

% Commande pour activer l'alternance en bleu
\newcommand{\monStyleTableau}{
  % On commence à la ligne 2 (pour laisser l'en-tête tranquille)
  % Alterne entre un bleu très clair et le blanc
  \rowcolors{2}{bleuTresClair}{white} 
}



\title{Rapport PFA}
\author{Islam ABOBAL}
\date{June 2026}

\usepackage{xcolor}
\usepackage[skins]{tcolorbox}

% Définition des couleurs personnalisées pour la page de garde
\definecolor{titleblue}{RGB}{23, 108, 172}
\definecolor{sectionblue}{RGB}{0, 162, 232}

\usepackage{booktabs}
\usepackage{tabularx}
\usepackage{array}
\usepackage{multirow}
\usepackage{longtable}
\usepackage{xcolor}
\usepackage{colortbl}

